\documentclass[a4paper,10pt]{report}\input{../0.Préambule/Préambule_cours_prof}\begin{document}

\pagestyle{empty}
\newpage
\noindent NOM, Prénom et classe : \dotfill

\section*{Les angles}
\addcontentsline{toc}{section}{Les angles}

\noindent \textit{L'usage de la calculatrice est \textbf{interdite}.} \hfill \textit{Durée : 30 minutes}

\medskip
\paragraph{Exercice 1}: \hfill \textbf{4,5 points}

\medskip
\noindent Observe les angles codés sur la figure ci dessous et complète la tableau :

\noindent 
\begin{minipage}{9cm}
\renewcommand{\arraystretch}{1.5}
\begin{tabularx}{\linewidth}{|m{5mm}|*{3}{>{\centering \arraybackslash}X|}}
\hline
n \degre & Notation & Sommet & Cotés \\\hline
\ding{192} & & & \\\hline
\ding{193} & & & \\\hline
\ding{194} & & & \\\hline
\end{tabularx}
\end{minipage}
\hspace{5mm}
\begin{minipage}{7cm}
\begin{center}
\includegraphics[scale=0.4]{../40.Image/DS6 Angles.png}
\end{center}
\end{minipage}

\medskip
\paragraph{Exercice 2}: \hfill \textbf{4 points}

\medskip
\noindent \begin{enumerate}
\item Trace un angle $\widehat{xAy}$ mesurant $32$\degre .
\item Trace un angle $\widehat{LAC}$ mesurant $107$\degre .
\end{enumerate}

\vspace{4cm}

\medskip
\paragraph{Exercice 3}: \hfill \textbf{6 points}

\medskip
\noindent Donne la mesure des angles ci-dessous.

\vspace{-5mm}
\begin{center}
\begin{tikzpicture}[line cap=round,line join=round,>=triangle 45,x=1.0cm,y=1.0cm,scale=0.55]
\clip(-3.,-2.) rectangle (28.,6.);
\draw [line width=1.2pt] (2.74,4.68)-- (-1.98,1.06);
\draw [line width=1.2pt] (-1.98,1.06)-- (7.,-1.);
\draw [line width=1.2pt] (10.8,4.98)-- (18.78,1.16);
\draw [line width=1.2pt] (18.78,1.16)-- (27.08,0.58);
\draw [line width=1.2pt] (6.2,3.94)-- (16.34,0.4);
\draw [line width=1.2pt] (6.2,3.94)-- (12.78,-1.28);
\begin{scriptsize}
\draw [color=black] (2.74,4.68)-- ++(-3.5pt,0 pt) -- ++(7.0pt,0 pt) ++(-3.5pt,-3.5pt) -- ++(0 pt,7.0pt);
\draw[color=black] (2.2,5.23) node {\large{$A$}};
\draw [color=black] (-1.98,1.06)-- ++(-3.5pt,0 pt) -- ++(7.0pt,0 pt) ++(-3.5pt,-3.5pt) -- ++(0 pt,7.0pt);
\draw[color=black] (-2.64,1.15) node {\large{$B$}};
\draw [color=black] (7.,-1.)-- ++(-3.5pt,0 pt) -- ++(7.0pt,0 pt) ++(-3.5pt,-3.5pt) -- ++(0 pt,7.0pt);
\draw[color=black] (7.42,-0.41) node {\large{$C$}};
\draw [color=black] (10.8,4.98)-- ++(-3.5pt,0 pt) -- ++(7.0pt,0 pt) ++(-3.5pt,-3.5pt) -- ++(0 pt,7.0pt);
\draw[color=black] (11.22,5.57) node {\large{$D$}};
\draw [color=black] (18.78,1.16)-- ++(-3.5pt,0 pt) -- ++(7.0pt,0 pt) ++(-3.5pt,-3.5pt) -- ++(0 pt,7.0pt);
\draw[color=black] (18.54,0.63) node {\large{$E$}};
\draw [color=black] (27.08,0.58)-- ++(-3.5pt,0 pt) -- ++(7.0pt,0 pt) ++(-3.5pt,-3.5pt) -- ++(0 pt,7.0pt);
\draw[color=black] (27.4,1.13) node {\large{$F$}};
\draw [color=black] (6.2,3.94)-- ++(-3.5pt,0 pt) -- ++(7.0pt,0 pt) ++(-3.5pt,-3.5pt) -- ++(0 pt,7.0pt);
\draw[color=black] (5.64,4.45) node {\large{$G$}};
\draw [color=black] (16.34,0.4)-- ++(-3.5pt,0 pt) -- ++(7.0pt,0 pt) ++(-3.5pt,-3.5pt) -- ++(0 pt,7.0pt);
\draw[color=black] (16.22,1.05) node {\large{$H$}};
\draw [color=black] (12.78,-1.28)-- ++(-3.5pt,0 pt) -- ++(7.0pt,0 pt) ++(-3.5pt,-3.5pt) -- ++(0 pt,7.0pt);
\draw[color=black] (12.04,-1.47) node {\large{$I$}};
\end{scriptsize}
\end{tikzpicture}
\end{center}

\vspace{-10mm}
\paragraph{Exercice 4}: \hfill \textbf{3 points}

\medskip
\noindent Pour chaque cas, donne la nature de l'angle (aigu, obtus, plat, droit)

\begin{enumerate}
\renewcommand{\arraystretch}{1.8}
\begin{tabularx}{\linewidth}{m{1cm}*{1}{>{\arraybackslash}X}m{0.5cm}m{1cm}*{1}{>{\arraybackslash}X}}
\item $12,3$\degre & \dotfill &&
\item $179,9$\degre & \dotfill \\
\item $90$\degre & \dotfill &&
\item $180$\degre & \dotfill \\
\item $154$\degre & \dotfill &&
\item $49$\degre & \dotfill \\
\end{tabularx}
\end{enumerate}

\paragraph{Exercice 5}: \hfill \textbf{2,5 points}

\medskip
\noindent \begin{minipage}{9cm}
Sur la figure ci-dessous, les mesures ne sont pas respectées ! 

\noindent $\widehat{ABC}$ est un angle plat. Déterminer la valeur de $\widehat{EBD}$.
\end{minipage}
\hspace{5mm}
\begin{minipage}{7cm}
\begin{center}
\begin{tikzpicture}[line cap=round,line join=round,>=triangle 45,x=1.0cm,y=1.0cm,scale =0.5]
\clip(10.5,0.2) rectangle (24.5,3.7);
\draw[line width=1.2pt,color=red,fill=red,fill opacity=0] (19,1) -- (19,2) -- (18.,2) -- (18.,1.) -- cycle; 
\draw [shift={(18.,1.)},line width=1.2pt,color=red,fill=red,fill opacity=0] (0,0) -- (161.56505117707798:2) arc (161.56505117707798:180.:2) -- cycle;
\draw [line width=1.2pt] (11.,1.)-- (18.,1.);
\draw [line width=1.2pt] (18.,1.)-- (24.,1.);
\draw [line width=1.2pt] (18.,1.)-- (12.,3.);
\draw [line width=1.2pt] (18.,1.)-- (18.,3.);
\begin{scriptsize}
\draw [color=black] (11.,1.)-- ++(-3.5pt,0 pt) -- ++(7.0pt,0 pt) ++(-3.5pt,-3.5pt) -- ++(0 pt,7.0pt);
\draw[color=black] (10.63,1.37) node {\large{$A$}};
\draw [color=black] (18.,1.)-- ++(-3.5pt,0 pt) -- ++(7.0pt,0 pt) ++(-3.5pt,-3.5pt) -- ++(0 pt,7.0pt);
\draw[color=black] (17.94,0.4) node {\large{$B$}};
\draw [color=black] (24.,1.)-- ++(-3.5pt,0 pt) -- ++(7.0pt,0 pt) ++(-3.5pt,-3.5pt) -- ++(0 pt,7.0pt);
\draw[color=black] (24.29,1.40) node {\large{$C$}};
\draw [color=black] (12.,3.)-- ++(-3.5pt,0 pt) -- ++(7.0pt,0 pt) ++(-3.5pt,-3.5pt) -- ++(0 pt,7.0pt);
\draw[color=black] (11.1,3.41) node {\large{$D$}};
\draw [color=black] (18.,3.)-- ++(-3.5pt,0 pt) -- ++(7.0pt,0 pt) ++(-3.5pt,-3.5pt) -- ++(0 pt,7.0pt);
\draw[color=black] (17.60,3.34) node {\large{$E$}};
%\draw[color=ffqqqq] (19.0,1.50) node {\large{$90\textrm{\degre}$}};
\draw[color=red] (15,1.5) node {\large{$27\textrm{\degre}$}};
\end{scriptsize}
\end{tikzpicture}
\end{center}
\end{minipage}


\end{document}